\section{Related Work}

\subsection{Deep Learning for ECG Classification}

Hannun \etal~\cite{hannun2019} reported cardiologist-level arrhythmia
detection from single-lead ambulatory recordings using a 34-layer
convolutional network trained on 53{,}549 patients; Ribeiro
\etal~\cite{ribeiro2020} extended this to twelve-lead ECGs over two million
recordings. On PTB-XL specifically, Strodthoff \etal~\cite{strodthoff2021}
established the benchmarking protocol used here and reported macro AUC of
0.925--0.937 for superdiagnostic classification with XResNet1D and
Inception1D backbones. Subsequent work has explored attention
mechanisms~\cite{yao2020}, self-supervised pretraining~\cite{kiyasseh2021},
transformers~\cite{yan2019,natarajan2020} and multi-task formulations.

We position our results relative to this literature honestly: the encoders
used here are standard and not tuned for peak accuracy, and their absolute
performance sits below the strongest published numbers. The object of study is
the accuracy--dimension relationship under nesting, which is measured within a
fixed training protocol and is therefore unaffected by that offset. Where
absolute performance matters---in the state-of-the-art comparison of
Section~\ref{sec:sota}---we say so explicitly rather than presenting adaptive
and non-adaptive systems as though they were competing on the same axis.

\subsection{Efficient Models for Edge Deployment}

Knowledge distillation~\cite{hinton2015} transfers a large teacher's output
distribution to a compact student, typically achieving 2--4$\times$
compression with modest accuracy loss~\cite{li2023kd}. Pruning~\cite{han2016}
removes redundant weights; quantisation~\cite{jacob2018} reduces numerical
precision. TinyML systems have demonstrated arrhythmia detection within
kilobyte-scale memory budgets~\cite{sakr2023,banbury2020}. Early-exit
networks~\cite{teerapittayanon2016} and slimmable
architectures~\cite{yu2019} allow a degree of runtime adaptivity, but through
specialised multi-branch or width-switchable designs.

Two distinctions are worth drawing carefully, since they are frequently
elided. First, these methods reduce \emph{inference compute}; MRL does not.
Second, they produce a distinct artefact per operating point; MRL does not.
The methods are therefore complementary rather than competing---quantising an
MRL encoder addresses the compute axis that nesting leaves untouched---and we
treat them as such throughout.

\subsection{Nested Learning and Matryoshka Representations}

Behrouz \etal~\cite{behrouz2025} formalised Nested Learning as a view of
machine learning models as systems of nested optimisation problems with
distinct context flows and update rates. Matryoshka Representation
Learning~\cite{kusupati2022} instantiates this in representation space: for a
predefined set of nesting dimensions, the first $m$ coordinates of the
embedding independently constitute a usable representation. On ImageNet this
yields up to 14$\times$ smaller embeddings at equivalent accuracy.
MatFormer~\cite{kudugunta2023} extends nesting to transformer weights, and
adaptive retrieval systems built on MRL are deployed at
scale~\cite{kusupati2023adanns}.

Within medicine, MRL has been applied to histopathology
segmentation~\cite{golts2025} and biomedical text
embeddings~\cite{wang2024pubmedbert}. Neither addresses time-series
biosignals, where the relevant structure---quasi-periodic morphology at
multiple timescales---differs substantially from both images and text.
Table~\ref{tab:related} positions this work.

\begin{table}[!t]
\centering
\caption{Positioning relative to existing approaches. ``Single artefact''
means one set of weights serves all operating points; ``adaptive dim.'' means
the operating point is selectable at inference without retraining;
``compute saving'' means inference cost falls with the operating point.}
\label{tab:related}
\small
\begin{tabular}{lcccc}
\toprule
Approach & Single & Adaptive & Compute & 1-D \\
         & artefact & dim. & saving & signals \\
\midrule
Distillation~\cite{hinton2015}      & No  & No      & Yes & Yes \\
Pruning / quant.~\cite{han2016,jacob2018} & No & No & Yes & Yes \\
TinyML ECG~\cite{sakr2023}          & No  & No      & Yes & Yes \\
Slimmable nets~\cite{yu2019}        & Yes & Partial & Yes & No  \\
Early exit~\cite{teerapittayanon2016} & Yes & Partial & Yes & Yes \\
MRL (vision/NLP)~\cite{kusupati2022} & Yes & Yes    & No  & No  \\
MatFormer~\cite{kudugunta2023}      & Yes & Yes     & Yes & No  \\
\midrule
This work                           & Yes & Yes     & No  & Yes \\
\bottomrule
\end{tabular}
\end{table}

Note that the ``compute saving'' column marks our own method \emph{No}. This
is the honest entry, and it is the column that the deployment argument in
Section~\ref{sec:compute} must therefore be built without.
